% ═══════════════════════════════════════════════════════════════════════════
\section{Lavori Correlati}\label{sec:lavori}
% ═══════════════════════════════════════════════════════════════════════════

Questo capitolo passa in rassegna la letteratura pertinente, organizzata per
temi: metodi di ottimizzazione stocastica, metodi con campionamento dinamico,
metodi di secondo ordine stocastici e regolarizzazione $L_1$.

\subsection{Metodi di ottimizzazione stocastica}

Il punto di partenza di tutti i metodi stocastici è il lavoro pionieristico di
Robbins e Monro~\cite{robbins1951}, che ha introdotto il metodo del gradiente
stocastico (SGD): ad ogni iterazione si stima il gradiente su un singolo
esempio (o un piccolo batch) e si aggiorna $w$ con un passo decrescente.
Il costo per iterazione è indipendente da $N$, ma la varianza dello stimatore
limita la velocità di convergenza e la precisione raggiungibile, come analizzato
in dettaglio da Bottou e Bousquet~\cite{bottou2008}, che hanno formalizzato il
compromesso tra costo per iterazione e velocità di convergenza su larga scala.

Per ridurre la varianza dello stimatore sono state proposte diverse famiglie di
metodi. I metodi \emph{a varianza ridotta} combinano la stima stocastica con
una correzione periodica basata su un gradiente (quasi) esatto: SVRG di Johnson
e Zhang~\cite{johnson2013} mantiene una copia dei parametri su cui ricalcola un
gradiente completo a intervalli regolari, mentre SAGA di Defazio
et~al.~\cite{defazio2014} mantiene una tavola dei gradienti individuali e
corregge la stima a ogni iterazione. Questi metodi ottengono una convergenza
lineare per problemi fortemente convessi, ma richiedono di memorizzare o
ricalcolare quantità globali e la dimensione del batch rimane comunque fissata
a priori. Una direzione diversa è rappresentata dai metodi adattivi, tra cui
Adam di Kingma e Ba~\cite{kingma2015}, che scala il passo per componente usando
stime dei momenti del gradiente; Adam è estremamente diffuso in pratica, ma le
sue proprietà di convergenza sono state a lungo oggetto di studio e la sua
analisi per problemi non convessi è delicata.

Tutti questi metodi condividono una caratteristica: la dimensione del batch è
\textbf{fissa} o comunque determinata da una regola prestabilita, non da una
condizione di accuratezza verificata durante l'esecuzione. È proprio questo il
punto che il campionamento dinamico intende superare.

\subsection{Campionamento dinamico}

L'idea di far crescere la dimensione del batch \emph{durante} l'ottimizzazione
non è nuova, ma è stata formalizzata in modo sistematico da Byrd, Chin, Nocedal
e Wu~\cite{byrd2012}, il lavoro di riferimento di questa tesi. Gli autori

Approcci correlati adottano regole di crescita della dimensione del batch
basate su criteri diversi. Ad esempio, una strategia semplice e diffusa è la
\emph{crescita geometrica} $n_k = \lceil a^k \rceil$, che però richiede la
conoscenza (o una stima) del numero di condizionamento $\kappa$. Alcuni metodi
propongono invece di adattare la dimensione del batch in modo euristico in base
all'andamento della funzione obiettivo o della perdita di validazione, senza
una garanzia teorica. Il vantaggio dell'approccio basato sulla varianza è
duplice: da un lato fornisce una \emph{giustificazione statistica} alla regola
di aggiornamento, dall'altro si adatta automaticamente alla difficoltà locale
del problema, senza richiedere la stima di $\kappa$.

\subsection{Metodi di secondo ordine stocastici}

I metodi del primo ordine convergono lentamente su problemi mal condizionati,
poiché il numero di condizionamento $\kappa$ del problema entra nel tasso di
convergenza. I metodi di Newton incorporano l'informazione di curvatura
dell'Hessiana e raggiungono un tasso di convergenza locale quadratico (o
superlineare), a costo di risolvere ad ogni iterazione un sistema lineare
$H_k d = -g_k$; si veda il testo di Nocedal e Wright~\cite{nocedal2006} per
una trattazione completa.

Su larga scala la costruzione esplicita dell'Hessiana è impraticabile, e sono
stati sviluppati metodi che ne usano soltanto \emph{approssimazioni}. I metodi
\emph{inexact Newton}~\cite{dembo1982} risolvono il sistema di Newton in modo
approssimato, controllando il residuo del risolutore interno tramite un
parametro di forcing; il metodo del gradiente coniugato (CG) applicato al
sistema di Newton, in forma \emph{Hessian-free}, permette di calcolare i
prodotti Hessiana-vettore senza costruire la matrice~\cite{shewchuk1994,
nocedal2006}. In ambito stocastico, Byrd et~al.~\cite{byrd2011} hanno proposto
un metodo che sottocampiona sia il gradiente sia l'Hessiana, mostrando che un
numero di iterazioni del CG che cresce con la precisione desiderata mantiene la
convergenza. Il metodo Newton-CG presentato in questa tesi si colloca in questa
linea di ricerca, con due innovazioni rispetto ai lavori precedenti: la
dimensione del campione per il gradiente è regolata dinamicamente dalla CCV, e
il criterio di arresto del CG interno è \emph{adattivo}, basato sulla varianza
campionaria dei prodotti Hessiana-vettore, in modo da non sprecare iterazioni
interne quando l'approssimazione dell'Hessiana è limitata dal sottocampionamento.

\subsection{Regolarizzazione $L_1$ in ottimizzazione}

La regolarizzazione $L_1$ produce soluzioni \emph{sparse} ed è uno strumento
fondamentale per l'interpretabilità dei modelli ad alta dimensionalità, a
partire dal LASSO di Tibshirani~\cite{tibshirani1996}. Il termine di penalità
$\nu\|w\|_1$ rende la funzione obiettivo non differenziabile, e i metodi classici
del primo ordine devono essere adattati. Una tecnica basilare è il
\emph{soft-thresholding}, usato nei metodi proximali (ad esempio ISTA e FISTA
di Beck e Teboulle~\cite{beck2009}), che applica alla soluzione l'operatore di
soglia $S_\nu(w_i) = \operatorname{sign}(w_i)\max\{|w_i| - \nu, 0\}$.

Per i problemi che coinvolgono anche un termine liscio $J(w)$, un'alternativa
efficace è rappresentata dai metodi a \emph{active set} e dalle strategie di
\emph{ottimizzazione su facce ortanti}: si identifica l'insieme delle variabili
che all'ottimo saranno nulle e si minimizza la funzione su tale faccia, dove il
termine $L_1$ è differenziabile. Questa è la strategia seguita, ad esempio, da
Byrd et~al.~\cite{byrd2011} per problemi regolarizzati e adottata anche nel
metodo Newton-CG-$L_1$ descritto in questa tesi, che combina l'identificazione
dell'active set con l'Hessiana sottocampionata e la ricerca lineare proiettata.

osservano che, per garantire che la direzione di ricerca sia una direzione di
discesa, è sufficiente che l'errore di approssimazione del gradiente sia
controllato in norma relativa dalla tolleranza $\theta$. Poiché l'errore non è
osservabile direttamente, essi lo stimano tramite la \emph{varianza campionaria}
dei gradienti individuali sul batch, ottenendo una condizione di controllo
della varianza (CCV): se la varianza stimata è troppo alta rispetto alla norma
del gradiente, la dimensione del batch viene aumentata secondo una formula
esplicita. L'analisi mostra che il numero totale di valutazioni di gradienti
per raggiungere una precisione $\varepsilon$ è $O(L/(\lambda\varepsilon))$,
indipendente da $N$ e competitivo con i migliori bound noti per SGD.
